% ========================================
% Name: Guilherme Armin Da Silva Anton
% UFID: 2641-9801
% COP3530 Project 2 - GeoPulse
% File: Report.tex
% Purpose: Source for the final GeoPulse Project 2 report.
% ========================================
\documentclass[12pt]{article}
\usepackage[margin=0.55in]{geometry}
\usepackage{array}
\usepackage{booktabs}
\usepackage{hyperref}
\usepackage{titlesec}
\usepackage{enumitem}
\hypersetup{colorlinks=true, urlcolor=blue, linkcolor=black}
\setlength{\parindent}{0pt}
\setlength{\parskip}{3pt}
\setlist{nosep,leftmargin=*}
\titlespacing*{\section}{0pt}{6pt}{3pt}
\titlespacing*{\subsection}{0pt}{4pt}{2pt}
\newcolumntype{L}[1]{>{\raggedright\arraybackslash}p{#1}}
\begin{document}

\begin{titlepage}
\centering
{\Large University of Florida COP3530 Project 2\par}
\vspace{0.4in}
{\Huge\bfseries GeoPulse\par}
\vspace{0.2in}
{\Large Spatial Index Benchmarking for NYC Collision Data\par}
\vfill
\begin{tabular}{ll}
Team Name: & GeoPulse Solo Team\\
Team Member: & Guilherme Armin Da Silva Anton\\
UFID: & 2641-9801\\
GitHub Username: & arminanton\\
Repository: & \url{https://github.com/arminanton/project-2}\\
Video Demo: & \url{https://youtu.be/REPLACE_WITH_FINAL_VIDEO}\\
Language: & C++14\\
Submission: & Final group deliverable report\\
\end{tabular}
\vfill
\textbf{Summary.} GeoPulse compares Linear Scan, KD-tree, and Quadtree
spatial search over the NYC Motor Vehicle Collisions dataset. The two
non-trivial comparable structures implemented from scratch are the KD-tree and
Quadtree; Linear Scan is retained as the correctness oracle and baseline.
\end{titlepage}

\section{Extended and Refined Proposal}
\textbf{Problem.} NYC collision data contains millions of records, but raw CSV
search does not make spatial questions easy. GeoPulse asks: how much faster do
spatial data structures answer location-based collision queries than a brute
force scan?\newline
\textbf{Motivation.} Transportation analysts, students, and residents may want
to ask questions such as ``what crash is closest to this location?'' or ``how
many crashes happened within half a mile?'' A spatial index can reduce repeated
query cost after one build step.\newline
\textbf{Features implemented.} GeoPulse loads and validates CSV data, stores
accepted collisions in memory, builds three indexes, runs nearest-neighbor,
radius, and rectangular range queries, verifies KD-tree and Quadtree answers
against Linear Scan, and prints a benchmark report. Recovery mode writes a
checkpoint and accepted-record cache so interrupted CSV loads can resume.\newline
\textbf{Data.} The public dataset is NYC Open Data's Motor Vehicle Collisions
Crashes table, export id \texttt{h9gi-nx95}. The uploaded snapshot contains
2,246,476 CSV rows and 29 columns. GeoPulse accepts 2,005,745 rows after strict
numeric latitude/longitude validation; 240,625 rows are missing coordinates and
106 contain non-empty invalid coordinates.\newline
\textbf{Tools.} The implementation uses C++14, CMake, Catch2, Git/GitHub,
Git LFS-ready data layout, VS Code and CLion project settings, optional
clang-format, clang-tidy, and cpplint. No library KD-tree or Quadtree is used.\newline
\textbf{Algorithms implemented.} Linear Scan checks every point and is used as
the oracle. KD-tree recursively chooses medians on alternating x/y axes. Quadtree
recursively partitions the bounding rectangle into four quadrants.\newline
\textbf{Additional structures.} Shared records/points vectors avoid duplicate
metadata in tree nodes. Bounding boxes support pruning. Result normalization
sorts ids before set comparison.\newline
\textbf{Role distribution.} This was completed as a solo project. I performed
proposal refinement, architecture, data ingestion, algorithms, testing,
benchmarking, documentation, and packaging.

\section{Architecture and Design Decisions}
The code follows single-responsibility modules: \texttt{core} owns geometry and
constants; \texttt{data} owns CSV parsing, validation, loading, and recovery;
\texttt{index} owns the three comparable spatial structures; \texttt{query}
owns coordinate-facing query helpers; \texttt{benchmark} owns timing and
correctness comparison; \texttt{app/ui} owns output and command dispatch.
\texttt{main.cpp} only creates \texttt{CliApplication} and delegates.

The major post-proposal refinement was adding a stricter loader. The proposal
counted rows with non-empty coordinates; the final code accepts only finite,
numeric, legal latitude/longitude values. I also added quoted-newline CSV
support after discovering that a physical-line parser can miscount this dataset.
Recovery was designed with an accepted-record cache because a byte offset alone
cannot restore records that were already parsed before a restart.

\begin{center}
\begin{tabular}{L{1.3in}L{2.1in}L{2.0in}}
\toprule
Structure & Representation & Reason \\
\midrule
Linear Scan & Non-owning pointer to all points & Simple oracle; every result is exact.\\
KD-tree & Unique-pointer binary nodes with axis field & Median partitions improve average repeated query time.\\
Quadtree & Unique-pointer quadrant nodes with capacity/depth limits & Spatial boxes give intuitive radius/range pruning.\\
\bottomrule
\end{tabular}
\normalsize
\end{center}

\section{Complexity Analysis}
Let $N$ be indexed points and $K$ be returned matches. Worst cases remain
linear when data is adversarial, duplicate-heavy, or the query covers most of the
map.

\begin{center}
\small
\begin{tabular}{L{2.1in}L{1.45in}L{2.85in}}
\toprule
Command / method & Worst-case & Explanation \\
\midrule
\texttt{--data / DataLoader::load} & $O(N)$ & Streams each CSV row once.\\
\texttt{LinearScan::build} & $O(1)$ & Stores a reference to the point vector.\\
\texttt{LinearScan::nearest} & $O(N)$ & Checks every accepted point.\\
\texttt{LinearScan::radiusQuery} & $O(N)$ & Checks every distance.\\
\texttt{LinearScan::rangeQuery} & $O(N)$ & Checks every bounding-box containment.\\
\texttt{KDTree::build} & $O(N\log N)$ avg, $O(N^2)$ worst & Median selection per level.\\
\texttt{KDTree::nearest} & $O(N)$ & Branch pruning may fail.\\
\texttt{KDTree::radiusQuery} & $O(N+K)$ & Large circles can touch all branches.\\
\texttt{KDTree::rangeQuery} & $O(N+K)$ & Large boxes can touch all branches.\\
\texttt{QuadTree::build} & $O(N\log N)$ avg, $O(ND)$ worst & $D$ is max depth.\\
\texttt{QuadTree::nearest} & $O(N)$ & Box pruning may fail.\\
\texttt{QuadTree::radiusQuery} & $O(N+K)$ & Large circles can touch all nodes.\\
\texttt{QuadTree::rangeQuery} & $O(N+K)$ & Large boxes can touch all nodes.\\
\texttt{BenchmarkRunner::run} & $O(N\log N+QN)$ & Includes baseline correctness checks.\\
\bottomrule
\end{tabular}
\end{center}

\section{Testing, Tooling, and Output}
Catch2 tests cover argument validation, exact public sample output, quoted
CSV fields, data-loader row counts, recovery cache reuse, geometry, empty-index
behavior, negative radius behavior, deterministic nearest ties, KD-tree and
Quadtree correctness against Linear Scan, benchmark correctness flags, and a
style guard for student files outside Catch2. The public
sample output is deterministic:
\begin{verbatim}
Dataset summary
total rows: 8
accepted coordinate rows: 6
excluded rows: 2
missing-coordinate rows: 1
invalid-coordinate rows: 1
\end{verbatim}

Local quality tooling includes CMake targets \texttt{Main} and \texttt{Tests},
configure-time sample data copying, GitHub settings, branch-safe hooks, Git LFS
attributes for snapshots, VS Code and CLion settings, \texttt{.clang-format}
for Allman braces and 80 columns, and \texttt{.clang-tidy} with analyzer,
bugprone, performance, readability, and modernize checks. Disabled tidy checks
are documented as style conflicts: trailing return types are not used in this
beginner-friendly C++14 code, auto is not forced, magic numbers are handled by
shared constants and tests, and identifier-length warnings conflict with common
loop/index names.

\section{Reflection and Lessons Learned}
Working solo meant all architecture, implementation, testing, and reporting had
to stay highly organized. The most important challenge was not writing the tree
code itself; it was making all structures comparable under the same distance
model and proving correctness against a baseline. Another challenge was the CSV
format: quoted newlines made it clear that robust data ingestion matters before
benchmarking can be trusted.

If I restarted the project, I would implement the public facade, exact sample
outputs, and style guard even earlier. I learned how to separate domain records
from lightweight index points, why recovery needs persisted accepted records,
how KD-tree and Quadtree pruning differ, and why benchmarks must verify answers
before interpreting speed.

\section{References}
NYC Open Data. \emph{Motor Vehicle Collisions - Crashes}, dataset id
\texttt{h9gi-nx95}.\newline
University of Florida COP3530 Project 2 assignment description and rubric.\newline
Catch2 unit testing framework documentation.

\end{document}
